\documentclass[10pt,a4paper]{amsart}
\usepackage[margin=19mm]{geometry}
\usepackage{amsmath,amssymb,mathtools}
\usepackage[scr=rsfs,cal=euler]{mathalfa}
\usepackage{xfrac}
\usepackage[unicode,hidelinks,bookmarks=false]{hyperref}
\newcommand{\ie}{i.e.\ }
\newcommand{\cf}{cf.\ }
\newcommand{\resp}{respectively\ }
\newcommand{\Subsection}{\subsection}
\providecommand{\nobreakdash}{\nobreak\hbox{-}}
\setlength{\emergencystretch}{2em}
\multlinegap0pt
\usepackage{amssymb}
\usepackage{bbm}
\usepackage[all]{xy}
\usepackage{tikz}
\newtheorem{thm}{Theorem}
\newtheorem{cor}{Corollary}
\newcommand{\F}{\mathbb F}
\title{Irreducible polynomials with restricted digits in base $b(T)$}
\author{Juan Ar\'evalo, Matilde Lal\'in}
\date{Source: arXiv:2609.25367v1 (CC BY 4.0)}
\begin{document}
\maketitle

\noindent\emph{Source and licence.} Irreducible polynomials with restricted digits in base $b(T)$. Version arXiv:2609.25367v1 (CC BY 4.0), distributed by arXiv under the Creative Commons Attribution 4.0 International licence, http://creativecommons.org/licenses/by/4.0/, as recorded in the arXiv licence metadata for this exact version and shown on the abstract page https://arxiv.org/abs/2609.25367v1. Full licence notice: \texttt{licenses/arxiv-2609.25367-CC-BY-4.0.txt}.\par\smallskip

\noindent\textit{Editorial scope.} Counted unit: the article's cor cor:cardinality (source0.tex lines 1551--1589). The article's complete original proof of it is delivered with it (source0.tex lines 1592--1664, one \texttt{proof} environment of 73 lines). No mathematics is written, repaired or re-derived: the statement and the proof are the article's own lines, in the article's own notation, and no retained line is substituted or re-flowed.  Retained uncounted as the source-local context the proof needs: lines 348--348; lines 583--590; lines 960--971.  Nothing was omitted from any retained span, so the package carries no omission record.  No retained line contains a citation, so the wrapper carries no bibliography and no bibliography entry.  No retained line contains a figure environment, an image-inclusion command or drawing code, so no image is delivered and the package declares no asset dependency.  Editorial formatting only, in addition: an AMS \texttt{amsart} preamble that restates the article's own declarations and macros used by the retained lines, namely the environments \texttt{thm}, \texttt{cor} on separate counters, together with the standard packages those lines need.  The retained environments print in this document as Theorem 1, Corollary 1 in order of appearance, whereas the article prints the same items with the numbering of its own full document.  The complete native source of the article as distributed in the arXiv e-print for this exact version is bundled unchanged as \texttt{sources/211294/source0.tex}.
\section{Setup and auxiliary results}\label{sec:setup}

\begin{equation}\label{eq:Parseval}
\sum_{r\neq0}
\left|\widehat{\mathbbm 1_A}(r)\right|^2
=\sum_{a,a'\in A} \sum_{r\not =0}\psi(\langle a-a',r\rangle)=
\sum_{a,a'\in A}\left(q^m\mathbbm{1}_{a=a'}-1\right)
=q^m|A|-|A|^2=
|A|(q^m-|A|),
\end{equation}

\begin{thm}\label{thm:main} Let $b\in\mathcal M_m$ and let $\varnothing\neq\mathcal R\subsetneq D$. Assume that \[ L_\infty(\mathcal R^c)<|\mathcal R^c|, \qquad \kappa(\mathcal R^c)<1. \]
If $\ell\geq4$, then \begin{equation*}%\label{eq:main} 
I_n^{\mathcal R} =\frac{|\mathcal M_n^{\mathcal R}|}{n}\Bigl(\mathfrak S(\mathcal R)+\mathcal E\Bigr), \end{equation*} where \[ |\mathcal M_n^{\mathcal R}|=q^\delta|\mathcal R^c|^{\ell-1} \qquad \mbox{ and }\qquad \mathfrak S(\mathcal R)=\frac{q^m}{\Phi(b)} \frac{|\mathcal R^c\cap D^\times|}{|\mathcal R^c|}.\] 
Moreover, \[ \mathcal E\ll(1+\log_q n)q^{E_\eta} +\left(nq^{\frac{1-\delta}{4}} +(1+\log_q n)q^{-\frac n4+\frac m2+\frac{3\delta}{4}}\right) \kappa(\mathcal R^c)^{\ell-1}, \] where, writing $\eta=\eta(\mathcal R^c)$, \[ E_\eta=\frac{m}{24}\left(4\eta+11-\sqrt{16\eta^2+(96\ell-56)\eta+49}\right). \] 
In particular, for fixed $q$, $b$, and $\mathcal R$, as $n\to\infty$, we have $\ell\rightarrow \infty$ and we get 
\[
I_n^{\mathcal R}
=
\frac{|\mathcal M_n^{\mathcal R}|}{n}
\left(\mathfrak S(\mathcal R)+o(1)\right).
\]
\end{thm}

\begin{cor}\label{cor:cardinality}
Let $b\in\mathcal M_m$, let $D=\F_q[T]/(b)$, and let
$\varnothing\neq\mathcal R\subsetneq D$ with $s=|\mathcal R|$.
Let $\ell$ and $\delta$ be defined as in Section~\ref{sec:setup},
and define
\[
A(q,m,s):=
\frac{q^m-s+\sqrt{(q^m-1)s(q^m-s)}}{q^m}.
\]
Suppose that $\ell\geq4$ and
\begin{equation}\label{eq:s-cond}
s<
\frac{q^{\frac{3m}{4}}
\left(q^{\frac m4}-1\right)}
{q^{\frac m2}+q^{\frac m4}+2}.
\end{equation}
Then
\[
I_n^{\mathcal R}
=
\frac{|\mathcal M_n^{\mathcal R}|}{n}
\left(\mathfrak S(\mathcal R)+\mathcal E\right),
\] where
\begin{align*}
\mathcal E\ll\;&
(1+\log_q n)
q^{\frac m{24}\left(13-\sqrt{25+48\ell}\right)}
+
\left(
nq^{\frac{1-\delta}{4}}
+
(1+\log_q n)
q^{-\frac n4+\frac m2+\frac{3\delta}{4}}
\right)
\left(
\frac{q^{\frac{3m}{4}}A(q,m,s)}{q^m-s}
\right)^{\ell-1}.
\end{align*}
\end{cor}

\begin{proof}
    

For $r\neq0$ we have $\widehat{\mathbbm 1_{D}}(r)=0$, and therefore
$\widehat{\mathbbm 1_{\mathcal R^c}}(r)
=
-\widehat{\mathbbm 1_{\mathcal R}}(r)$. Since $\left|\widehat{\mathbbm 1_{\mathcal R}}(r)\right|\leq s,$ we conclude that $\left|\widehat{\mathbbm 1_{\mathcal R^c}}(r)\right|\leq s$ and
\[
L_\infty(\mathcal R^c)\leq s.
\]
We have $\widehat{\mathbbm{1}_{D}}(0)=q^m$, which implies $\widehat{\mathbbm{1}_{\mathcal{R}^c}}(0)=q^m- \widehat{\mathbbm{1}_{\mathcal{R}}}(0)= q^m-s$. Thus
\[
\sum_{r\in \F_q^m}|\widehat{\mathbbm{1}_{\mathcal{R}^c}}(r)|=(q^m-s)+\sum_{r\ne0}|\widehat{\mathbbm{1}_{\mathcal{R}}}(r)|.
\]

By Parseval \eqref{eq:Parseval},
\[
\sum_{r\not =0}|\widehat{\mathbbm{1}_{\mathcal{R}}}(r)|^2
=s(q^m -s).
\]

By Cauchy--Schwarz
\[
\left(\sum_{r\not =0}|\widehat{\mathbbm{1}_{\mathcal{R}}}(r)|\right)^2\leq
\left(\sum_{r\not=0}1\right)\left(\sum_{r\not =0}|\widehat{\mathbbm{1}_{\mathcal{R}}}(r)|^2\right)
=(q^m-1)\sum_{r\not=0}|\widehat{\mathbbm{1}_{\mathcal{R}}}(r)|^2=(q^{m}-1)s(q^m-s),
\]
so $\sum_{r\not=0}|\widehat{\mathbbm{1}_{\mathcal{R}}}(r)|\leq \sqrt{(q^m-1)s (q^m-s)}$ and
\[
L_1(\mathcal R^c)\leq A(q,m,s).
\]
The assumption \eqref{eq:s-cond} on $s$ is equivalent to
\[
\frac{q^{\frac{3m}{4}}A(q,m,s)}{q^m-s}<1.
\]
Since $L_1(\mathcal R^c)\leq A(q,m,s)$, it follows that
\[
\kappa(\mathcal R^c)
=
\frac{q^{\frac{3m}{4}}L_1(\mathcal R^c)}{q^m-s}
<1.
\]
Moreover, by \eqref{eq:s-cond}, we have 
\[s<\frac{q^{\frac{3m}{4}}
\left(q^{\frac m4}-1\right)}
{q^{\frac m2}+q^{\frac m4}+2}<\frac{q^m}{q^\frac{m}{2}+1}\]
and this implies
\[
\frac{s}{q^m-s}<q^{-\frac{m}{2}}.
\]
Since $L_\infty(\mathcal R^c)\leq s$, it follows that
\[
\frac{L_\infty(\mathcal R^c)}{|\mathcal R^c|}
=
\frac{L_\infty(\mathcal R^c)}{q^m-s}
<q^{-\frac{m}{2}}<1.
\]
Thus $L_\infty(\mathcal R^c)<|\mathcal R^c|$, and
\[
\eta(\mathcal R^c)
=
\frac1m\log_q
\frac{|\mathcal R^c|}{L_\infty(\mathcal R^c)}
>\frac12.\]
Since $E_\eta$ is decreasing as a function of $\eta$ for $\ell\geq4$, we have
\[
E_\eta\leq E_{1/2}
=
\frac m{24}\left(13-\sqrt{25+48\ell}\right).
\]
The result now follows from Theorem~\ref{thm:main}.

\end{proof}

\end{document}
